% ===================================================================
%  TABEL Nr. 5 — Maja ja selle ehitamine.
%  Traduction 2026 d'apres texto/io, controlee sur texto/fr et
%  texto/fi. Consigne : voir texto/et/10-tabel-01.tex.
%
%  LE TABLEAU 5 ROMANCHE AVAIT PARTAGE CE CHANTIER EN TROIS COUCHES ;
%  L'ESTONIEN LE PARTAGE EN DEUX, ET LA COUPURE TOMBE AU MEME ENDROIT.
%
%  Le romanche avait la pierre latine — « murader », « cuvrider »,
%  « faver » —, l'atelier allemand — « schrinari », de Schreiner —,
%  et le papier roman du dehors. On en avait tire que la ligne passe
%  entre ce qui se fait a la main, ce qui se fait a l'etabli et ce
%  qui se fait sur le papier.
%
%  L'ESTONIEN N'A PAS DE LATIN, ET IL ARRIVE AU MEME RESULTAT PAR UN
%  TOUT AUTRE CHEMIN. Il a un mot native, « sepp », le forgeron, et
%  il en a fait le nom general de l'artisan :
%    « müürsepp »   le macon, le forgeron-de-mur ;
%    « puusepp »    le charpentier, le forgeron-de-bois ;
%    « klaassepp »  le vitrier ;
%    « lukksepp »   le serrurier ;
%    « kingsepp »   le cordonnier ;
%    « kullassepp » l'orfevre.
%  La ou l'allemand a six mots sans rapport — Maurer, Zimmermann,
%  Glaser, Schlosser, Schuster, Goldschmied — l'estonien en derive
%  six d'un seul.
%
%  ET UN SEUL METIER A RESISTE : LE MENUISIER, QUI SE DIT « tisler »,
%  DE Tischler. C'est exactement celui que le romanche avait pris a
%  l'allemand sous la forme « schrinari », de Schreiner. Deux langues
%  sans contact, deux mecanismes opposes — l'une herite du latin,
%  l'autre derive d'une racine a elle — et le meme trou au meme
%  endroit.
%
%  LA CAUSE N'EST DONC DANS AUCUNE DES DEUX LANGUES, ELLE EST DANS LE
%  METIER. Un macon travaille la pierre du lieu avec une truelle ; un
%  menuisier travaille avec un ETABLI, des rabots, un tour et un
%  apprentissage regle — c'est-a-dire avec un atelier, et l'atelier
%  est arrive du nord avec son nom. C'est la deuxieme fois que la
%  serie verifie une regle sur deux langues sans rapport : le tableau
%  11 romanche l'avait fait sur la caisse d'epargne et les pompiers.
%
%  Comme dans texto/io, aucun alinea de ce tableau n'est numerote par
%  l'auteur : la cle est « 01 » pour tout le tableau.
%
%  UNE INVERSION, AU BLOC t05-01-49, ET C'EST LA MEME CAUSE QU'AU
%  TABLEAU 3 : la chaine de genitifs estonienne place le possesseur
%  devant, « trepi kolm astet » et « välisukse lävi », la ou l'ido
%  range marches (14), perron (9), seuil (10), porte (8). On brise la
%  chaine en deux, comme au tableau 3 : les marches d'abord, le
%  perron en apposition, puis le seuil avec une relative pour la
%  porte. Le procede est le meme et il resservira.
% ===================================================================

%%K t05-apar-1 apar
\VUsaut{6.0mm}
\VUtitre{15.0pt}{60}{TABEL Nr. 5}
\VUsaut{1.4mm}
\VUtitre{11.4pt}{0}{Maja ja selle ehitamine.}
\VUsaut{2.2mm}

%%K t05-tit-1 sub
\VUpk{10.2pt}{40}{Hea kohtumine.}
\VUsaut{3.0mm}

%%K t05-01-1 p
\textsc{Jaan}. --- Onu, ma tahaksin väga teada, kuidas ehitatakse \VUgras{maja}.

%%K t05-01-2 p
\textsc{Onu} (80). --- See on väga lihtne, mu sõber, tule minuga, ma näitan sulle
seda, mida härra Bernard (79) laseb ehitada ja kus ma hakkan elama.

%%K t05-01-3 p
\textsc{Jaan}. --- Kes siis joonistas selle maja \VUgras{plaani} \textsuperscript{(76)}?

%%K t05-01-4 p
\textsc{Onu}. --- \VUgras{Arhitekt} \textsuperscript{(75)}; ta esitas väga selge
joonise, mis kiideti heaks, ja \VUgras{ettevõtja} \textsuperscript{(77)} alustas töid.

%%K t05-01-5 p
\textsc{Jaan}. --- Kes siis on ettevõtja?

%%K t05-01-6 p
\textsc{Onu}. --- See on härra Valge, sinu sõbra Jaakobi isa. Ta juhib töölisi koos
härra Punasega, arhitektiga.

%%K t05-01-7 p
Noh! Vaat just õigel ajal tulevad härra Valge oma \VUgras{joonlauaga}
\textsuperscript{(78)} ja härra Punane oma plaaniga.

%%K t05-01-8 p
Tere päevast, härrad. Kuidas teil läheb?

%%K t05-01-9 p
\textsc{Härra Punane}. --- Väga hästi, aitäh. Tere päevast, Jaan, kui suureks see laps
on kasvanud!

%%K t05-01-10 p
\textsc{Onu}. --- Jah, tõesti, ta on väike mees. Ta on väga tõsine, talle tuleb
seletada kõike, mida ta näeb. Täna tahaks ta teada, kuidas ehitatakse maja.

%%K t05-tit-2 sub
\VUpk{10.2pt}{40}{Töölised.}
\VUsaut{3.0mm}

%%K t05-01-11 p
\textsc{Härra Valge}. --- Noh, tule minuga, sõbrake, ma seletan sulle seda kohe, sest
ma armastan väga poisse, kes tahavad õppida.

%%K t05-01-12 p
Sa näed seda töölist, kes hoiab käes \VUgras{labidat} \textsuperscript{(82)}: ta on
\VUgras{kaevaja} \textsuperscript{(81)}. Ta kaevab \VUgras{kraavi} \textsuperscript{(46)},
et paigaldada veetorustik. Ta on kaevanud ka pinnase seal, kus maja tuleb, et müürsepp
saaks püstitada neli \VUgras{müüri}. Nende müüride osa, mis on maa sees, kannab nime
\VUgras{vundament} \textsuperscript{(2)}. Vundamentide vahele jääv ruum moodustab
\VUgras{keldri} \textsuperscript{(3)}. Sellel keldril on väike aken, mida nimetatakse
\VUgras{keldriaknaks} \textsuperscript{(5)}, \VUgras{uks} \textsuperscript{(4)} ja trepp.

%%K t05-01-13 p
\VUgras{Müürsepp} \textsuperscript{(108)} jätkas müüride ehitamist, jättes avad
\VUgras{ustele} \textsuperscript{(8)} ja \VUgras{akendele} \textsuperscript{(17)}.

%%K t05-01-14 p
\textsc{Jaan}. --- Aga seda müürseppa ma ei näe!

%%K t05-01-15 p
H\textsuperscript{ra} \textsc{Valge}. --- Nüüd on ta majas töö lõpetanud. Ta on
\VUgras{talli} \textsuperscript{(54)} juures, oma \VUgras{tellingutel}
\textsuperscript{(110)}, ta ehitab \VUgras{pesuköögi} \textsuperscript{(56)} müüri.

%%K t05-01-16 p
Tal on kellu, küna ja \VUgras{lood} \textsuperscript{(109)}. Aga sa ei jäta neid nimesid
meelde.

%%K t05-01-17 p
\textsc{Jaan}. --- Kindlasti jätan, sest ma palun kohe oma onult väikest kellu ja küna,
ja ma teen ise nöörist loodi.

%%K t05-tit-3 sub
\VUsaut{3.0mm}
\VUpk{10.2pt}{40}{Materjal.}
\VUsaut{3.0mm}

%%K t05-01-18 p
\textsc{Onu}. --- Ah! väga hästi. Aga ütle mulle, Jaan, mida on vaja maja ehitamiseks?

%%K t05-01-19 p
\textsc{Jaan}. --- Vaja on \VUgras{tahutud kive} \textsuperscript{(59)} ja
\VUgras{mörti} \textsuperscript{(65)}.

%%K t05-01-20 p
\textsc{Onu}. --- Õige, vaata seda meest suure kübaraga oma \VUgras{vankril}
\textsuperscript{(136)}. Ta on \VUgras{kutsar} \textsuperscript{(135)}. Iga päev toob ta
siia \VUgras{kiviplokke} \textsuperscript{(60)}. Ta laadib oma vankri õues maha, ja
\VUgras{kivitahujad} \textsuperscript{(83)} tahuvad plokke ja saevad neid oma
\VUgras{kivisaega} \textsuperscript{(85)}. \VUgras{Käsitööline} \textsuperscript{(146)}
viib need müürsepale, kes paneb need paika.

%%K t05-01-21 p
Vahepeal valmistab \VUgras{mördisegaja} \textsuperscript{(87)} mörti. Tal on vaja
\VUgras{liiva} \textsuperscript{(62)}, \VUgras{lupja} \textsuperscript{(63)} ja
\VUgras{vett} \textsuperscript{(64)}. Lubi on tema kõrval \VUgras{kotis}
\textsuperscript{(71)}, \VUgras{müürsepa abiline} \textsuperscript{(111)} toob talle liiva
\VUgras{käruga} \textsuperscript{(112)} ja \VUgras{õpipoiss} \textsuperscript{(113)} on pumba
(58) juures. Ta täidab \VUgras{ämbri} \textsuperscript{(114)}. Mört on peagi valmis.
\VUgras{Mördikandja} \textsuperscript{(107)} võtab selle ja viib mööda redelit üles
tellingutele, kus töötab müürsepp, kes lõpetab seda maja külge.

%%K t05-01-22 p
Ja nüüd, kuidas nimetatakse töölist, kes teeb \VUgras{katust} \textsuperscript{(31)}?

%%K t05-01-23 p
\textsc{Jaan}. --- Ta on \VUgras{katusetegija} \textsuperscript{(118)}.

%%K t05-tit-4 sub
\VUsaut{3.0mm}
\VUpk{10.2pt}{40}{Maja katus.}
\VUsaut{3.0mm}

%%K t05-01-24 p
H\textsuperscript{ra} \textsc{Valge}. --- Jah, aga kõigepealt tuleb \VUgras{puusepp}
\textsuperscript{(115)}.

%%K t05-01-25 p
Puusepp teeb \VUgras{puitkonstruktsiooni} \textsuperscript{(35)}. Ta ühendab
\VUgras{talad} \textsuperscript{(37)} \VUgras{tappidega} \textsuperscript{(117)}. Need
töölised seal üleval, oma laiade sametpükstega, on puusepad. Üks hoiab väikest tala ja
teine raiub tapi oma \VUgras{kirvega} \textsuperscript{(116)}. See, kes on
\VUgras{korstna} \textsuperscript{(41)} juures, on katusetegija. Ta lööb latid (*)
palkidele ja \VUgras{kiltkivid} \textsuperscript{(66)} lattidele. Mõnikord paneb ta
katusekive.

%%K t05-noto-1 noto
\VUnotes{143.2mm}{%
(*) \VUgras{Balk-o} = S. \textit{Balken}, I. \textit{joist}, P. \textit{solive}, It.
\textit{travicello}, V. \textit{balka}, Hisp. Port. \textit{viga}. Tehniline tüvi, mida
seni ei ole vastu võetud, kuid mida tuleks pakkuda prantsuse sõna \textit{solive}
mõtte edasiandmiseks; see ei ole ei tala P. \textit{poutre} ega väike tala. See on
nelinurkseks tahutud puitlatt, mis kannab põrandat ja lage.%
}

%%K t05-01-26 p
Talli katus on \VUgras{katusekividega kaetud} \textsuperscript{(55)}, maja oma on
\VUgras{kiltkiviga kaetud} \textsuperscript{(32)} ja veranda oma on \VUgras{tsingist}
\textsuperscript{(24)}.

%%K t05-01-27 p
\textsc{Jaan}. --- Kas katusetegija teeb ka tsinkkatuseid?

%%K t05-01-28 p
H\textsuperscript{ra} \textsc{Valge}. --- Ei, vaid \VUgras{tsingitööline}
\textsuperscript{(121)} või \VUgras{pliitööline}, see, kes paneb maja katusele
\VUgras{katuseakna} \textsuperscript{(38)}. Ta paneb ka \VUgras{vihmaveetorud}
\textsuperscript{(43)} ja \VUgras{räästarennid} \textsuperscript{(42)}. Ta joodab kokku
tsingi ja pleki. Ta kinnitas \VUgras{piksevarda} \textsuperscript{(40)} ja
\VUgras{tuulelipu} \textsuperscript{(39)} \VUgras{viilu} \textsuperscript{(36)} otsa.

%%K t05-01-29 p
\textsc{Jaan}. --- Kas maja on siis valmis?

%%K t05-tit-5 sub
\VUsaut{3.0mm}
\VUpk{10.2pt}{40}{Tööriistad.}
\VUsaut{3.0mm}

%%K t05-01-30 p
H\textsuperscript{ra} \textsc{Valge}. --- Mitte päris. \VUgras{Krohvija}
\textsuperscript{(99)} ehitab \VUgras{tellistest} \textsuperscript{(61)} vaheseinad, mis
jaotavad selle eri tubadeks. Ta krohvib \VUgras{kipsiga} \textsuperscript{(70)}
\VUgras{laed} \textsuperscript{(26)} ja müürid. Nüüd teeb ta \VUgras{veranda}
\textsuperscript{(33)} lage. Tal on, nagu müürsepal, \VUgras{kellu} \textsuperscript{(100)}
ja \VUgras{küna} \textsuperscript{(101)}.

%%K t05-01-31 p
Siis teeb tisler uksed, aknad, \VUgras{parketi} \textsuperscript{(16)} ja
\VUgras{aknaluugid} \textsuperscript{(19)}.

%%K t05-01-32 p
Nüüd töötab ta õues oma \VUgras{tööpingi} \textsuperscript{(90)} juures. Tal on
\VUgras{saag} \textsuperscript{(94)} puidu (69) saagimiseks, \VUgras{höövel}
\textsuperscript{(91)}, \VUgras{kruustangid} \textsuperscript{(92)}, \VUgras{peitel}
\textsuperscript{(94 bis)}, \VUgras{sirkel} \textsuperscript{(95)} ja puust
\VUgras{haamer} \textsuperscript{(95 bis)}.

%%K t05-01-33 p
\textsc{Jaan}. --- Ah! aga puutööd teha on lõbusam kui müüri laduda. Onu, sa ostad
mulle tööpingi, sae ja hööveli, sest ma teen varsti Medorile kuudi.

%%K t05-01-34 p
\textsc{Onu}. --- Jah, poisike.

%%K t05-01-35 p
\textsc{Jaan}. --- Kui rahul ma olen! Ja kes see on, kes seisab välisukse juures?

%%K t05-tit-6 sub
\VUsaut{3.0mm}
\VUpk{10.2pt}{40}{Värvimine.}
\VUsaut{3.0mm}

%%K t05-01-36 p
H\textsuperscript{ra} \textsc{Valge}. --- Ta on \VUgras{maaler} \textsuperscript{(102)}.
Ta seisab oma redelil \VUgras{värvi} \textsuperscript{(103)} potiga ja oma
\VUgras{pintsliga} \textsuperscript{(104)}. Ta värvib välisukse puitu, et see paremini
säiliks.

%%K t05-01-37 p
Ta liimib ka korterites tapeete.

%%K t05-01-38 p
\VUgras{Tapeetija} \textsuperscript{(107)}, kes on \VUgras{redelil}
\textsuperscript{(106)}, \VUgras{rõdul} \textsuperscript{(28)}, paneb \VUgras{rulood}
\textsuperscript{(29)}.

%%K t05-01-39 p
Tööline, kes seisab suure akna juures, on \VUgras{klaassepp} \textsuperscript{(96)}. Ta
kinnitab \VUgras{aknaklaasid} \textsuperscript{(18)} \VUgras{kitiga} \textsuperscript{(73)}.
See teine tööline, kes põlvitab keldriukse juures, on \VUgras{lukksepp}
\textsuperscript{(97)}. Ta paneb \VUgras{lukku} \textsuperscript{(6)}. Tema \VUgras{viil}
\textsuperscript{(98)} on tema kõrval.

%%K t05-01-40 p
\textsc{Jaan}. --- Kas on ka lukksepp see tööline, kes lööb oma \VUgras{haamriga}
\textsuperscript{(84)} rauale?

%%K t05-01-41 p
H\textsuperscript{ra} \textsc{Valge}. --- Täpsemalt öeldes on ta sepp (125). Ta sepistab
\VUgras{rauast osi} \textsuperscript{(67)} \VUgras{elektrikule} \textsuperscript{(124)}, kes
on \VUgras{vankrikuuri} \textsuperscript{(51)} katusel. Tal on \VUgras{sepikoda}
\textsuperscript{(126)}, \VUgras{alasi} \textsuperscript{(127)} ja \VUgras{tangid}
\textsuperscript{(128)}.

%%K t05-01-42 p
Elektrik kinnitab telefoni \VUgras{vasktraadi} \textsuperscript{(44)}.

%%K t05-tit-7 sub
\VUpk{10.2pt}{40}{Aiapidamine.}
\VUsaut{3.0mm}

%%K t05-01-43 p
\textsc{Onu}. --- \VUgras{Õue} \textsuperscript{(45)} tagaosas, \VUgras{võre}
\textsuperscript{(47)} ja \VUgras{värava} \textsuperscript{(48)} juures näed sa
\VUgras{aednikku} \textsuperscript{(129)}. Ta valmistab ette \VUgras{aiakese}
\textsuperscript{(49)}. Ta on kaevanud maad oma \VUgras{labidaga} \textsuperscript{(131)},
ta külvab seemneid ja riisub \VUgras{rehaga} \textsuperscript{(130)}. Siis kastab ta oma
\VUgras{kastekannuga} \textsuperscript{(132)} \VUgras{peenrad} \textsuperscript{(150)} ja
taimed hakkavad kasvama.

%%K t05-01-44 p
\VUgras{Tallipoiss} \textsuperscript{(133)}, kes on just hobust hooldanud ja talle süüa
andnud, puhastab \VUgras{tõlda} \textsuperscript{(52)} käsna, \VUgras{käsipumba}
\textsuperscript{(134)} ja veeämbriga. Siis viib ta selle tagasi vankrikuuri ja sulgeb
ukse paksu \VUgras{riiviga} \textsuperscript{(53)}.

%%K t05-01-45 p
Nüüd oled sa rahul, ma arvan, sa said küllalt seletusi.

%%K t05-01-46 p
\textsc{Jaan}. --- Jah, onu, aga ma näeksin meelsasti maja sisemust.

%%K t05-01-47 p
\textsc{Onu}. --- See juhtub homme. Härra Punane seletab sulle plaani ja näitab sulle
iga toa nime.

%%K t05-tit-8 sub
\VUsaut{3.0mm}
\VUpk{10.2pt}{40}{Maja sisemine korraldus.}
\VUsaut{2.4mm}

%%K t05-apar-2 apar
{\centering\textit{(Vaata plaani.)}\par}
\VUsaut{2.4mm}

%%K t05-01-48 p
\textsc{Arhitekt}. --- Tule, Jaan, ma lasen sul kohe uurida maja eri osi.

%%K t05-01-49 p
Astume \VUgras{esimesele korrusele} \textsuperscript{(7)}. Siin on \VUgras{uksekella}
\textsuperscript{(11)} nupp. Me astume üles kolm \VUgras{astet} \textsuperscript{(14)}
--- see on \VUgras{välistrepp} \textsuperscript{(9)} ---, me ületame \VUgras{läve}
\textsuperscript{(10)}, mis on \VUgras{välisukse} \textsuperscript{(8)} ees, ja vaat
et me olemegi \VUgras{trepi} \textsuperscript{(13)} jalamil.

%%K t05-01-50 p
\textsc{Jaan}. --- See on koridor.

%%K t05-01-51 p
\textsc{Arhitekt}. --- Ei, see on \VUgras{esik} \textsuperscript{(12)}.
\VUgras{Koridor} \textsuperscript{(a)} on otsas. Paremal on \VUgras{salong}
\textsuperscript{(b)} ja \VUgras{söögituba} \textsuperscript{(c)}, söögitoa vastas on
\VUgras{köök} \textsuperscript{(d)} ja \VUgras{teenijatetuba} \textsuperscript{(e)}, siis
ees \VUgras{kabinet}\textsuperscript{(f)}. \VUgras{Veranda} \textsuperscript{(23)} oma
\VUgras{klaasuksega} \textsuperscript{(25)} on salongi kõrval.

%%K t05-01-52 p
Me läheme kohe trepist üles. Hoia end \VUgras{käsipuust} \textsuperscript{(15)} ja loe
astmeid. Neid on neliteist. Siin on \VUgras{trepimade} \textsuperscript{(m)}, me oleme
teisel \VUgras{korrusel} \textsuperscript{(27)}.

%%K t05-tit-9 sub
\VUsaut{3.0mm}
\VUpk{10.2pt}{40}{Teine korrus.}
\VUsaut{3.0mm}

%%K t05-01-53 p
Trepi vastas on \VUgras{käimla} \textsuperscript{(20)} ja \VUgras{tualettruum}
\textsuperscript{(j)}. Paremal ilus \VUgras{magamistuba} \textsuperscript{(g)} kahe
aknaga maja \VUgras{fassaadi} \textsuperscript{(1)} poole. Vasakul on \VUgras{vannituba}
\textsuperscript{(1)} ja \VUgras{külalistetuba} \textsuperscript{(i)} suure
\VUgras{alkooviga} \textsuperscript{(h)}: magamistoa kõrval on \VUgras{lastetuba}
\textsuperscript{(k)}. See avaneb avarale \VUgras{terrassile} \textsuperscript{(21)}.
Selle terrassi all on teine käimla (20) ja seina küljes on \VUgras{kell}
\textsuperscript{(22)}, mis kutsub õhtusöögile.

%%K t05-01-54 p
\textsc{Jaan}. --- Lähme \VUgras{kolmandale korrusele}.

%%K t05-01-55 p
\textsc{Arhitekt}. --- Kolmandat korrust ei ole. Katuse all on ainult
\VUgras{mansardid} \textsuperscript{(33)} ja pööning (34). Me oleme uurimise lõpetanud,
me võime tagasi alla minna.

%%K t05-01-56 p
\textsc{Jaan}. --- Aitäh, härra, see on väga huvitav. Ma jutustan kohe oma isale kõik,
mida ma nägin.

\VUsaut{4.0mm}
\VUfilet{22mm}
